\begin{proofbox}
\textbf{证明前统一约定}: $P(x_0,y_0)$ 在
$\frac{x^2}{a^2}+\frac{y^2}{b^2}=1$ 上, $a>b>0$, 且 $x_0y_0\ne0$.
以下六种方法均要得到 $k_lk_{OP}=-b^2/a^2$.

\textbf{方法一: 最快证明——直接使用切线方程 (推荐考场首选)}

\begin{enumerate}
  \item 椭圆在 $P$ 处的切线为
  \[
  \frac{x_0x}{a^2}+\frac{y_0y}{b^2}=1.
  \]
  移项并除以 $y_0/b^2$:
  \[
  y=-\frac{b^2x_0}{a^2y_0}x+\frac{b^2}{y_0},
  \qquad k_l=-\frac{b^2x_0}{a^2y_0}.
  \]
  \item 过 $O(0,0)$、$P(x_0,y_0)$ 的直线斜率是
  $k_{OP}=\frac{y_0}{x_0}$.
  \item 两式相乘, $x_0,y_0$ 直接约掉:
  \[
  \boxed{k_lk_{OP}=
  \left(-\frac{b^2x_0}{a^2y_0}\right)
  \frac{y_0}{x_0}=-\frac{b^2}{a^2}.}
  \]
  \item 因为 $c^2=a^2-b^2$, $e=c/a$, 故
  \[
  e^2-1=\frac{a^2-b^2}{a^2}-1=-\frac{b^2}{a^2}.
  \]
\end{enumerate}
\textbf{最快记忆}: \textbf{切线斜率含 $x_0/y_0$; 中心连线斜率含 $y_0/x_0$;
一乘只剩 $-b^2/a^2$.}

\textbf{方法二: 完全不用导数——切点重根 / 判别式思想}

\begin{enumerate}
  \item 设过 $P$ 的待定斜率直线为
  \[
  y-y_0=k(x-x_0).
  \]
  令 $t=x-x_0$, 则直线上任意点为
  $(x,y)=(x_0+t,y_0+kt)$.
  \item 代入椭圆方程:
  \[
  \frac{(x_0+t)^2}{a^2}+
  \frac{(y_0+kt)^2}{b^2}=1.
  \]
  展开并使用 $\frac{x_0^2}{a^2}+\frac{y_0^2}{b^2}=1$, 得
  \[
  t^2\left(\frac1{a^2}+\frac{k^2}{b^2}\right)
  +2t\left(\frac{x_0}{a^2}+\frac{ky_0}{b^2}\right)=0.
  \]
  \item $t=0$ 对应已知交点 $P$. 直线与椭圆相切意味着
  $P$ 是\textbf{二重交点}, 即 $t=0$ 为二重根.
  由于 $\frac1{a^2}+\frac{k^2}{b^2}>0$, 必须有
  \[
  \frac{x_0}{a^2}+\frac{ky_0}{b^2}=0.
  \]
  \item 解出
  \[
  k=k_l=-\frac{b^2x_0}{a^2y_0}.
  \]
  再乘 $k_{OP}=y_0/x_0$, 得
  \[
  \boxed{k_lk_{OP}=-\frac{b^2}{a^2}.}
  \]
\end{enumerate}
\textbf{方法价值}: 不必先记切线方程; “已知交点 + 相切 $\Rightarrow$ 重根”
是圆锥曲线直线相切问题的通用技巧.

\textbf{方法三: 隐函数求导——直接求切线斜率}

\begin{enumerate}
  \item 对椭圆方程两边关于 $x$ 求导:
  \[
  \frac{2x}{a^2}+\frac{2y}{b^2}\frac{dy}{dx}=0.
  \]
  \item 解得
  \[
  \frac{dy}{dx}=-\frac{b^2x}{a^2y}.
  \]
  \item 在 $P(x_0,y_0)$ 处,
  \[
  k_l=\left.\frac{dy}{dx}\right|_P
  =-\frac{b^2x_0}{a^2y_0}.
  \]
  \item 由 $k_{OP}=y_0/x_0$, 立即得到
  \[
  \boxed{k_lk_{OP}=-\frac{b^2}{a^2}.}
  \]
\end{enumerate}
\textbf{方法价值}: 会求隐函数导数时很快, 还能用于其他隐式曲线的切线题.

\textbf{方法四: 参数法——用一个角同时表示点和切线}

\begin{enumerate}
  \item 设
  \[
  x_0=a\cos\theta,\qquad y_0=b\sin\theta,
  \qquad \sin\theta\cos\theta\ne0.
  \]
  因为 $\cos^2\theta+\sin^2\theta=1$, 该点一定在椭圆上.
  \item 中心连线斜率为
  \[
  k_{OP}=\frac{b\sin\theta}{a\cos\theta}.
  \]
  \item 对参数方程求导:
  \[
  \frac{dx}{d\theta}=-a\sin\theta,\qquad
  \frac{dy}{d\theta}=b\cos\theta.
  \]
  因而切线斜率为
  \[
  k_l=\frac{dy/d\theta}{dx/d\theta}
  =-\frac{b\cos\theta}{a\sin\theta}.
  \]
  \item 相乘:
  \[
  k_lk_{OP}=-\frac{b\cos\theta}{a\sin\theta}
   \frac{b\sin\theta}{a\cos\theta}
   =\boxed{-\frac{b^2}{a^2}}.
  \]
\end{enumerate}
\textbf{方法价值}: 参数 $\theta$ 全部约去, 直观看出乘积与切点位置无关.

\textbf{方法五: 几何本质——把单位圆拉伸成椭圆}

\begin{enumerate}
  \item 从单位圆 $u^2+v^2=1$ 出发. 在非轴点 $Q$ 处,
  圆的半径 $OQ$ 与切线垂直, 因而斜率 $m_r,m_t$ 满足
  \[
  m_rm_t=-1.
  \]
  \item 对平面作横向放大 $a$ 倍、纵向放大 $b$ 倍的变换:
  \[
  (u,v)\longmapsto(x,y)=(au,bv).
  \]
  变换后的曲线为 $x^2/a^2+y^2/b^2=1$;
  $OQ$ 变为 $OP$, 原圆的切线变为椭圆在 $P$ 处的切线.
  \item 设原直线方向变化量为 $(\Delta u,\Delta v)$,
  则变换后为 $(a\Delta u,b\Delta v)$, 其斜率从
  $m=\Delta v/\Delta u$ 变成
  \[
  k=\frac{b\Delta v}{a\Delta u}=\frac ba m.
  \]
  故 $k_{OP}=\frac ba m_r$, $k_l=\frac ba m_t$.
  \item 两式相乘:
  \[
  \boxed{k_lk_{OP}=\left(\frac ba\right)^2m_tm_r
  =-\frac{b^2}{a^2}.}
  \]
\end{enumerate}
\textbf{方法价值}: 这是最直观的\textbf{几何解释}.
注意: 非等比拉伸\textbf{不保持角度}, 所以圆上的垂直关系通常不会在椭圆中保留下来.

\textbf{方法六: 中点弦法——先证明更广的结论, 再让弦趋于切线}

\begin{enumerate}
  \item 取椭圆上不同两点 $A(x_1,y_1)$、$B(x_2,y_2)$,
  令弦 $AB$ 的中点为 $M(u,v)$, 即
  \[
  u=\frac{x_1+x_2}{2},\qquad v=\frac{y_1+y_2}{2}.
  \]
  暂设 $x_1\ne x_2$、$u\ne0$, 以保证下面的斜率存在.
  \item 分别写出 $A,B$ 的椭圆方程并相减:
  \[
  \frac{x_1^2-x_2^2}{a^2}
  +\frac{y_1^2-y_2^2}{b^2}=0.
  \]
  用平方差公式拆开:
  \[
  \frac{(x_1-x_2)(x_1+x_2)}{a^2}
  +\frac{(y_1-y_2)(y_1+y_2)}{b^2}=0.
  \]
  \item 代入 $x_1+x_2=2u$, $y_1+y_2=2v$, 两边除以
  $2(x_1-x_2)$:
  \[
  \frac{u}{a^2}
   +\frac{v}{b^2}\frac{y_1-y_2}{x_1-x_2}=0.
  \]
  令 $k_{AB}=\frac{y_1-y_2}{x_1-x_2}$, 得
  \[
  k_{AB}\frac vu=-\frac{b^2}{a^2}.
  \]
  当 $M\ne O$ 且斜率都存在时, $v/u=k_{OM}$, 因而
  \[
  \boxed{k_{AB}k_{OM}=-\frac{b^2}{a^2}.}
  \]
  \item 让 $A,B$ 沿椭圆同时趋于同一非轴点 $P$.
  弦 $AB$ 的方向趋于 $P$ 处的切线方向, 中点 $M$ 趋于 $P$;
  因此 $k_{AB}\to k_l$, $k_{OM}\to k_{OP}$.
  取极限得到
  \[
  \boxed{k_lk_{OP}=-\frac{b^2}{a^2}.}
  \]
\end{enumerate}
\textbf{方法价值}: 把\textbf{弦的中点定理}与\textbf{切线}联系起来;
无需微积分就可先获得平行弦中点的定值性质, 再由割线趋于切线得到本题结论.
弦经过中心时 $M=O$, 不可写 $k_{OM}$; 轴方向的情况也须单独讨论.
\end{proofbox}
